\documentclass[../../main.tex]{subfiles}
\begin{document}

$n$を正の整数とし，
$\displaystyle \left( \cos\dfrac{2\pi}{n}k,\sin\dfrac{2\pi}{n}k \right)$
を座標とする点を$Q_n$であらわす．
このとき，$n$個の点$Q_0,Q_1,\cdots\cdots,Q_{n-1}$によって円周$x^2+y^2=1$は$n$等分される．

平面上の点$P$の座標を$(a,b)$とし，
$\displaystyle s_n=\dfrac{1}{n}({\overline{PQ_0}}^2+\overline{PQ_1}^2+\cdots\cdots+\overline{PQ_{n-1}}^2)$
とするとき，
$\displaystyle{\lambda}_P=\lim_{n\to\infty}s_n$の値を$a$，$b$をもちいてあらわせ．
また，$P$がどこにあれば${\lambda}_P$の値は最小となるか．

\end{document}
